
\documentclass{article}
\usepackage{colortbl}
\usepackage{makecell}
\usepackage{multirow}
\usepackage{supertabular}

\begin{document}

\newcounter{utterance}

\centering \large Interaction Transcript for game `eqbench', experiment `eq\_bench', episode 106 with reflexion{-}Qwen3.5{-}27B.
\vspace{24pt}

{ \footnotesize  \setcounter{utterance}{1}
\setlength{\tabcolsep}{0pt}
\begin{supertabular}{c@{$\;$}|p{.15\linewidth}@{}p{.15\linewidth}p{.15\linewidth}p{.15\linewidth}p{.15\linewidth}p{.15\linewidth}}
    \# & \multicolumn{2}{c}{Player} && \multicolumn{2}{c}{Game Master} \\
    \hline

    \theutterance \stepcounter{utterance}  
    & & & \multicolumn{4}{p{0.6\linewidth}}{
        \cellcolor[rgb]{0.9,0.9,0.9}{
            \makecell[{{p{\linewidth}}}]{
                \texttt{\tiny{[P1$\langle$GM]}}
                \texttt{You have attempted this same task before and did not succeed. These are the plans you wrote after those attempts.} \\
\texttt{Trial \#1: New plan: My previous attempt failed because I assigned a non{-}zero score to "Dismissive" (2) while ignoring the strong textual evidence for "Frustrated" as the primary emotion, likely leading to an incorrect ranking or intensity distribution. The dialogue shows Terry escalating from defensive ("If I have to") to angry ("You're goddamn right") and finally to a stern, frustrated ultimatum ("That's enough," "liability I can't afford"). Terry is not dismissive; he is deeply engaged and agitated by Cass's provocation. Next time, I will strictly analyze the emotional trajectory: Terry's repeated interruptions and defensive justifications indicate high frustration. I will assign a high score (e.g., 9 or 10) to Frustrated, and strictly 0 to Relieved, Dismissive, and Content, as the text provides no evidence of relief, dismissal, or contentment. I will ensure the scores reflect the singular dominant emotion driven by the conflict.} \\
\\ 
\\ 
\\ 
\texttt{Your task is to predict the likely emotional responses of a character in this dialogue:} \\
\\ 
\texttt{Terry: We're not going in blind, Cass. Not this time.} \\
\texttt{Cass: And who's gonna stop me? You?} \\
\texttt{Terry: If I have to. } \\
\texttt{Cass: Oh, right. The big, bad Terry Holloway. Always with a plan. } \\
\texttt{Terry: You're goddamn right. You know why? Because plans save lives.} \\
\texttt{Cass: Or maybe they just make you feel better about yourself. Like you're in control.} \\
\texttt{Terry: This isn't about control, Cass. It's about survival. } \\
\texttt{Cass: Survival? Or just your guilt?} \\
\texttt{Terry: That's enough.} \\
\texttt{Cass: Hit a nerve, did I? } \\
\texttt{Terry: You're good, Cass. But you're reckless. And that's a liability I can't afford.} \\
\texttt{Cass: And there it is. The threat. You gonna bench me, Terry?} \\
\texttt{Terry: If it means keeping you alive, yes.} \\
\texttt{Cass: I don't need your protection. I need to do my job.} \\
\texttt{{[}End dialogue{]}} \\
\\ 
\texttt{At the end of this dialogue, Terry would feel...} \\
\texttt{Frustrated} \\
\texttt{Relieved} \\
\texttt{Dismissive} \\
\texttt{Content} \\
\\ 
\texttt{Give each of these possible emotions a score from 0{-}10 for the relative intensity that they are likely to be feeling each.} \\
\\ 
\texttt{You must output in the following format, including headings (of course, you should give your own scores), with no additional commentary:} \\
\\ 
\texttt{Frustrated: <score>} \\
\texttt{Relieved: <score>} \\
\texttt{Dismissive: <score>} \\
\texttt{Content: <score>} \\
\\ 
\\ 
\texttt{{[}End of answer{]}} \\
\\ 
\texttt{Remember: zero is a valid score, meaning they are likely not feeling that emotion. You must score at least one emotion > 0.} \\
\\ 
\texttt{Your answer:} \\
            }
        }
    }
     \\ \\

    \theutterance \stepcounter{utterance}  
    & \multicolumn{4}{p{0.6\linewidth}}{
        \cellcolor[rgb]{0.561,0.737,0.561}{
            \makecell[{{p{\linewidth}}}]{
                \texttt{\tiny{[P1$\rangle$GM]}}
                \texttt{Frustrated: 10} \\
\texttt{Relieved: 0} \\
\texttt{Dismissive: 0} \\
\texttt{Content: 0} \\
            }
        }
    }
    & & \\ \\

    \theutterance \stepcounter{utterance}  
    & & & \multicolumn{2}{p{0.3\linewidth}}{
        \cellcolor[rgb]{0.9,0.9,0.9}{
            \makecell[{{p{\linewidth}}}]{
                \texttt{\tiny{[GM$|$GM]}}
                \texttt{Frustrated: 10} \\
\texttt{Relieved: 0} \\
\texttt{Dismissive: 0} \\
\texttt{Content: 0} \\
            }
        }
    }
    & & \\ \\

    \theutterance \stepcounter{utterance}  
    & & & \multicolumn{2}{p{0.3\linewidth}}{
        \cellcolor[rgb]{0.9,0.9,0.9}{
            \makecell[{{p{\linewidth}}}]{
                \texttt{\tiny{[GM$|$GM]}}
                \texttt{Frustrated: 8} \\
\texttt{Relieved: 0} \\
\texttt{Dismissive: 5} \\
\texttt{Content: 0} \\
            }
        }
    }
    & & \\ \\

    \theutterance \stepcounter{utterance}  
    & & & \multicolumn{2}{p{0.3\linewidth}}{
        \cellcolor[rgb]{0.9,0.9,0.9}{
            \makecell[{{p{\linewidth}}}]{
                \texttt{\tiny{[GM$|$GM]}}
                \texttt{game\_result = LOSE} \\
            }
        }
    }
    & & \\ \\

    \theutterance \stepcounter{utterance}  
    & & \multicolumn{4}{p{0.6\linewidth}}{
        \cellcolor[rgb]{0.95,0.95,0.95}{
            \makecell[{{p{\linewidth}}}]{
                \texttt{\tiny{[P1$\leftrightarrow$P]}}
                \texttt{Post{-}episode Reflexion (not sent to GM):} \\
\texttt{New plan: My previous attempt failed because I assigned a score of 0 to "Dismissive," overlooking the specific textual evidence where Terry explicitly devalues Cass's input and agency. While "Frustrated" is the dominant emotion, Terry's statement "You're good, Cass. But you're reckless" and his decision to bench her ("If it means keeping you alive, yes") are acts of dismissal regarding her capability and role. By assigning 0 to Dismissive, I likely failed to capture the nuance required by the target state, which probably expects a non{-}zero score for this secondary but present emotion. Next time, I will analyze the dialogue for both the primary emotional state (Frustrated) and the functional emotional response (Dismissive). I will assign a high score (e.g., 9 or 10) to Frustrated due to the conflict and a moderate score (e.g., 6 or 7) to Dismissive to reflect Terry's rejection of Cass's autonomy, while keeping Relieved and Content at 0.} \\
            }
        }
    }
    & \\ \\

\end{supertabular}
}

\end{document}
